\documentclass[conference]{IEEEtran}
\usepackage{cite}
\usepackage{amsmath}
\usepackage[hidelinks]{hyperref}

% Introduction for the Defense 4 DNP3 CLRT timing-obfuscation paper.
% Rewritten 2026-08-07 to the accepted bounded claim after the Phase 6 gate closed with
% TIMING EXPERIMENTS PARTIAL WITH CLOSED CLAIM BOUNDARY (see
% defense4/timing/evidence/EXPERIMENTAL_EVIDENCE_FREEZE.md and defense4/PAPER-STATE.md).
% Story order: (1) the timing fingerprint threat; (2) why general traffic obfuscation cannot work for
% ICS/SCADA, stated plainly; (3) the specific requirements that follow; (4) our in-network approach;
% (5) contributions bounded strictly to the accepted evidence. The thesis is stated directly, not held
% back. Claims about the not-yet-run controlled negatives and the cross-device classification are
% explicitly NOT made. No em dashes; plain, direct IEEE prose.

\title{An In-Network Timing Framework for Concealing the\\Response-Time Fingerprint of a Separate-Acknowledgment DNP3 Device}
\author{\IEEEauthorblockN{Anonymous submission}}

\begin{document}
\maketitle

\section{Introduction}

An attacker studies a control system before it acts on it. Fingerprinting is how a passive observer
turns ordinary network traffic into knowledge about the devices that produce it, learning a device's
type, vendor, and role without decoding any payload~\cite{kohno2005remote, radhakrishnan2014gtid,
shu2006fingerprint}. In industrial control systems the payoff is large and durable, because the field
devices are long lived and their roles are fixed, and prior work has fingerprinted supervisory and
field equipment directly from its traffic~\cite{jeon2016passive, formby2016control}. One of the
features that gives a device away is timing. A DNP3 outstation answers a master's poll with a
transport acknowledgment and then, a short and characteristic delay later, the response itself. That
delay, the cross-layer response time, distinguishes control-system devices by vendor and
model~\cite{formby2016control}. Because DNP3 carries it in the clear, an on-path observer reads the
fingerprint from routine polling traffic.

The obvious answer is traffic obfuscation, and there is a large toolkit for it: add cover packets, pad
sizes, and reshape flows so the observer's measurements no longer separate one source from
another~\cite{wright2009morphing, cai2014csbuflo, juarez2016wtfpad, apthorpe2019stp}. That toolkit does
not transfer to industrial control, and it is worth stating plainly why, because the reason shapes the
rest of the paper. General obfuscation was built for the encrypted Internet, where the traffic is
opaque and the shaper is free to invent packets and reshape flows because the observer cannot tell a
real exchange from a manufactured one~\cite{juarez2016wtfpad, wright2009morphing}, and where the
weaker of these countermeasures have in any case been broken by stronger
analysis~\cite{dyer2012peekaboo, sirinam2018df}. Legacy industrial traffic inverts every assumption
those methods rest on:

\begin{itemize}
\item \textbf{It is plaintext and checksummed.} A DNP3 frame is readable and each link block carries a
validating CRC~\cite{east2009taxonomy, fovino2010modbus}, so an injected dummy or an altered byte is
detectable by inspection, not hidden by encryption. Cover traffic and byte edits, the staples of
general obfuscation, are off the table.
\item \textbf{The endpoints cannot be changed.} The outstation is vendor firmware and DNP3 is a
deployed standard~\cite{east2009taxonomy, cardenas2011attacks}, so a defense cannot ask the device or
the protocol to shape their own traffic.
\item \textbf{Correctness is not optional.} Every poll must receive its real response, in order, within
the protocol's timing, retransmission, and quality-of-service bounds. A defense cannot drop, reorder,
or delay past those bounds without breaking the control loop.
\item \textbf{The fingerprint is in the timing.} Size- and volume-oriented defenses do not touch the
response-time signature that identifies the device.
\end{itemize}

The consequence is a specific and different problem. A defense for this setting must live in the
network, because the device and the protocol are fixed; it must preserve every original byte, because
the traffic is plaintext and checksummed; it must preserve the exchange and the correctness of every
delivered reply; and it must act on the timing observable itself, which is the thing that leaks. None
of the general obfuscation methods meets all four constraints at once, and that gap, not a missing
tuning of an existing method, is what this paper addresses.

We present an in-network defense that meets the four constraints. One programmable switch sits between
the DNP3 master and the field device and changes only when the device's acknowledgment and response
become externally visible, never their contents. Rather than editing any packet, it holds the original
acknowledgment and response in its queues while small internal marker packets recirculate to keep their
place, and the traffic manager releases the held packets on a schedule the control plane
sets~\cite{wang2020pinot, meier2022ditto, kfoury2021p4survey}. This design works with, rather than
against, the fact that a data-plane pipeline cannot pause an enqueued packet and recall it later at a
software-chosen time. Configurable gates on the acknowledgment and on the response give a small family
of timing transformations, from releasing the acknowledgment only when its matching response is ready,
to holding the acknowledgment or the response to a fixed deadline, to holding both, all as modes of one
mechanism selected at run time. We assume the switch is a trusted observation and control point on the
plaintext path, and that the adversary is a passive observer on the master-facing segment, where only
the released, shaped timing is visible.

\noindent\textbf{Contributions.} Each contribution is bounded by the evidence we report.
\begin{itemize}
\item We state why general traffic obfuscation does not work for legacy ICS traffic, and we frame the
cross-layer response time, identified by Formby et al.~\cite{formby2016control} as a device-
fingerprinting feature, as a reachable in-network shaping target under four constraints, in-network,
byte-preserving, correctness-preserving, and timing-targeted, that separate this problem from the
size- and encryption-oriented defenses above.
\item We design a unified, queue-resident timing framework whose event, acknowledgment-deadline,
response-deadline, and dual-deadline transformations are configurations of one mechanism. A
separate-acknowledgment device releases its acknowledgment before its response, so the response
obligation must survive the acknowledgment release; we show how the transaction state is kept alive
across that release so a later response is still held rather than forwarded unshaped.
\item We resolve the DNP3, TCP, and Tofino constraints the design meets: matching each transaction to
its DNP3 transaction-generation code on plaintext traffic, keeping the internal markers off the
observable wire, preserving the original bytes, and fitting the arithmetic within a single Tofino-1
ingress pipeline.
\item We implement the framework on an Intel Tofino-1 switch and evaluate it on silicon against a
physical SEL-751 relay. On the corrected implementation, the response-deadline and dual-deadline modes
normalize the measured response time to a fixed value and hold every response, verified by a
fail-closed measurement pipeline and reproduced across two campaigns; we quantify the resulting
collapse of the timing observable. We report these results, and their limits, for one relay, for READ
polling transactions, and for the response-time observable only. We do not claim device
indistinguishability across vendors, size concealment, or the full set of controlled failure cases;
the controlled negative-path tests and a cross-device classification study are stated as future work,
not as results.
\end{itemize}

\bibliographystyle{IEEEtran}
\bibliography{../../defense3/references}

\end{document}
